\section{Proof of the lifting lemma used in \texorpdfstring{\cref{lem:T0}}{Lemma 7.2}}\label{app:lifting}
The tiling lemma (\cref{lem:T0}) uses the standard statement that the lower facets of a lifted polytope project onto its shadow with disjoint interiors (the regular subdivision induced by a lifting, \cite{Ziegler1995,DLRS2010}; the construction behind the Bohne--Dress correspondence \cite{Bohne1992,RGZ1994,HRS2000,BLSWZ1999}). For completeness the argument is given here.

Let $\pi:\R^{d+1}\to\R^d$ forget the last coordinate and let $\tilde P\subset\R^{d+1}$ be a polytope with non-empty interior, $P=\pi(\tilde P)$. A face of $\tilde P$ is \emph{lower} if it is the set of maximisers of some linear functional $\nu=(\nu',\nu_{d+1})$ with $\nu_{d+1}<0$ on $\tilde P$; a lower facet is a lower face of dimension $d$.
\begin{lemma}[lifting]\label{lem:lift}
(i) On every lower face $G$ the projection $\pi$ is an affine isomorphism onto $\pi(G)$; (ii) the sets $\pi(F)$, $F$ a lower facet, cover $P$; (iii) the interiors of $\pi(F)$ and $\pi(F')$ are disjoint for distinct lower facets $F\ne F'$.
\end{lemma}
\begin{proof}
Put $h(x)=\min\{t:(x,t)\in\tilde P\}$ for $x\in P$.

(i) Let $G$ be the set of maximisers of $\nu$ with $\nu_{d+1}<0$ and $(x,t)\in G$. If $(x,t')\in\tilde P$ with $t'<t$, then $\nu\cdot(x,t')-\nu\cdot(x,t)=\nu_{d+1}(t'-t)>0$, contradicting maximality; hence $t=h(x)$. So $G\subseteq\{(x,h(x))\}$ and $\pi|_G$ is injective, and an injective affine map on a polytope is an affine isomorphism onto its image; in particular $\dim\pi(G)=\dim G\le d$.

(ii) The set $\pi(\mathrm{int}\,\tilde P)$ is open (since $\pi$ is an open map), convex, and dense in $P$; for a convex set $C$ with $\overline C=P$ one has $\mathrm{int}\,P=\mathrm{int}\,C$, so $\pi(\mathrm{int}\,\tilde P)=\mathrm{int}\,P$. Let $x\in\mathrm{int}\,P$ and $y=(x,h(x))$. The fibre $\tilde P\cap(\{x\}\times\R)$ contains a point $(x,t_0)$ of $\mathrm{int}\,\tilde P$, and $y\in\partial\tilde P$ by minimality of $h(x)$. Let $\nu$ be a supporting functional at $y$, $G_\nu=\{z\in\tilde P:\nu\cdot z=\nu\cdot y\}$. Since $(x,t)\in\tilde P$ for $t\ge h(x)$ near $t_0$ we have $\nu_{d+1}\le0$; and $\nu_{d+1}=0$ would give $\nu\cdot(x,t_0)=\nu\cdot y$, i.e.\ an interior point of $\tilde P$ on a supporting hyperplane, which is impossible. Hence $\nu_{d+1}<0$, so $y$ lies in a lower face $G_\nu$ and $x\in\pi(G_\nu)$. Thus $\mathrm{int}\,P\subseteq\bigcup\pi(G)$ over lower faces $G$. The lower faces of dimension $<d$ project (by (i)) to finitely many polytopes of dimension $<d$, a closed set with empty interior; a lower face of dimension $d$ is a lower facet. So the union of $\pi(F)$ over lower facets is a closed set that contains a dense subset of $\mathrm{int}\,P$, hence contains $\mathrm{int}\,P$ and, being closed, $P=\overline{\mathrm{int}\,P}$.

(iii) Let $F\ne F'$ be lower facets and $x$ an interior point of both $\pi(F)$ and $\pi(F')$. By (i) the points of $F,F'$ above $x$ are both equal to $y=(x,h(x))$, so $y\in F\cap F'$; as $\pi|_F$, $\pi|_{F'}$ are affine isomorphisms onto $d$-dimensional sets and $x$ is interior to the images, $y$ lies in the relative interior of both $F$ and $F'$. But a point of a polytope lies in the relative interior of exactly one face, so $F=F'$. \qedhere
\end{proof}
\begin{remark}
For the zonotope $\tilde Z=\sum[0,1]\tilde v_i$ of \cref{lem:T0}, full-dimensionality holds because, for $n>d$ and generic heights, some lifted generator is not in the span of the $d$ others of a basis; the lemma then gives that the lower facets $\sum_{i\in A_S}\tilde v_i+Z(\tilde v_S)$ project onto tiles $T_S$ that cover $Z(V)$ with disjoint interiors, which is the claim of \cref{lem:T0}.
\end{remark}
